% Generated draft: requires mathematical and source review.
\documentclass[11pt]{article}
\usepackage[margin=25mm]{geometry}
\usepackage{fontspec}
\setmainfont{Latin Modern Roman}
\usepackage{amsmath,amssymb,mathtools,unicode-math}
\setmathfont{Latin Modern Math}
\usepackage{enumitem,needspace,xurl,hyperref}
\setlength{\parindent}{0pt}
\setlength{\parskip}{5pt}
\setlength{\emergencystretch}{3em}
\newcommand{\N}{\mathbb{N}}
\newcommand{\Z}{\mathbb{Z}}
\newcommand{\Q}{\mathbb{Q}}
\newcommand{\R}{\mathbb{R}}
\newcommand{\C}{\mathbb{C}}
\newcommand{\F}{\mathbb{F}}
\newcommand{\A}{\mathbb{A}}
\newcommand{\PP}{\mathbb{P}}
\newcommand{\E}{\mathbb{E}}
\newcommand{\eps}{\varepsilon}
\newcommand{\dd}{\,\mathrm{d}}
\begin{document}
\section*{UnsolvedMath ID: 10400068}
\textbf{Source problem code:} catalogue entry
\subsection{Definitions and mathematical statement}
Let \(T\) be a two-strand tangle in the source's tangle category, and let \(Z(T)\) be its Kontsevich invariant, with the source's framing, endpoint, coefficient and completion conventions. Write \(K_T=\mathcal C(T)\) for the knot constructed in Figure 12: take the two-parallel of \(T\) and add the particular solid-line tangle shown in that figure. The dotted lines in the figure allow strands to be knotted and linked. Let \(P^\theta_{K_T}\) denote the source's two-loop polynomial of this knot. The requested construction is an explicit presentation of \(P^\theta_{\mathcal C(T)}\) in terms of \(Z(T)\), valid for every admissible \(T\). In schematic form, construct and justify a rule \(\Phi\) with \[ P^\theta_{\mathcal C(T)}=\Phi(Z(T)). \] Here \(\mathcal C\) and \(\Phi\) are descriptive placeholders, not new definitions of the missing tangle or of the desired formula. Figure 12 is not reproduced in the supplied text, so its crossing, connection and parallel conventions cannot be recovered from the caption alone. The excerpt also does not define the variables, relations or normalization of \(P^\theta\); consequently an explicit algebraic target must retain those source-dependent conventions rather than silently choose a two-loop invariant.
\subsection{Short English statement}
Kricker asks for a formula for the two-loop polynomial of a knot built from a two-strand tangle, expressed through that tangle's Kontsevich invariant. The exact knot construction depends on the absent Figure 12.
\subsection{Sources}
\begin{itemize}
\item Local UnsolvedMath registry, permanent ID \(10400068\), exact source code \texttt{catalogue entry}.
\item A. Kricker, attribution in Problem 3.18 of Ohtsuki, \emph{Problems on invariants of knots and 3-manifolds} (2002), PDF page 67 and Figure 12: \url{https://msp.org/gtm/2002/04/gtm-2002-04-024s.pdf}. Registry extraction: pdf-font-structure; recorded access date 2026-07-29.
\end{itemize}
\end{document}
